% !TeX root = ../algebra-1.tex
\chapter{Eigenvalues, diagonalisation and determinants}
\label{chap:algebra-eigenvalues}

\section{Eigenvectors and eigenspaces}

\begin{definition}
Let $A$ be an $n\times n$ matrix over $F$.
A \emph{nonzero} vector $v\in F^n$ is an \emph{eigenvector} with
\emph{eigenvalue} $\lambda\in F$ if
\[
 Av=\lambda v.
\]
The \emph{eigenspace} for $\lambda$ is
\[
 E_\lambda(A)=N(A-\lambda I_n).
\]
\end{definition}

An eigenvector is a direction preserved by the map: its image is a
scalar multiple of itself. Zero is never an eigenvector, although zero
can be an eigenvalue.
The eigenspace is a subspace and includes zero; its nonzero vectors are
exactly the eigenvectors for that eigenvalue.

\begin{example}[The all-ones matrix]
Let $J$ be the $3\times3$ matrix all of whose entries are $1$.
For $v=(a,b,c)^{\mathsf T}$,
$Jv=(a+b+c)(1,1,1)^{\mathsf T}$.
It follows that
\begin{align*}
 E_0(J)&=\{(a,b,c)^{\mathsf T}:a+b+c=0\}\\
 &=\spanop\bigl((-1,1,0)^{\mathsf T},(-1,0,1)^{\mathsf T}\bigr),\\
 E_3(J)&=\spanop\bigl((1,1,1)^{\mathsf T}\bigr).
\end{align*}
These are the only eigenvalues: if $Jv=\lambda v$ and $\lambda\ne0$,
then $v$ is a multiple of $(1,1,1)^{\mathsf T}$, which forces $\lambda=3$.
\end{example}

\begin{example}[The field matters]
For $A=\begin{pmatrix}0&1\\-1&0\end{pmatrix}$, we have $A^2=-I_2$.
If $Av=\lambda v$ with $v\ne0$, then $\lambda^2=-1$, so there are
no real eigenvalues or real eigenvectors.
Over $\mathbb C$ there are two eigenvalues:
\[
 A\begin{pmatrix}-i\\1\end{pmatrix}
 =i\begin{pmatrix}-i\\1\end{pmatrix},\qquad
 A\begin{pmatrix}i\\1\end{pmatrix}
 =-i\begin{pmatrix}i\\1\end{pmatrix}.
\]
\end{example}

\section{When is a matrix diagonalisable?}

\begin{definition}
$A$ is \emph{diagonalisable over $F$} if there is an invertible matrix
$P$ over $F$ such that $P^{-1}AP$ is diagonal.
\end{definition}

\begin{theorem}[Eigenvector basis criterion]
$A$ is diagonalisable over $F$ if and only if $F^n$ has a basis of
eigenvectors of $A$.
If $P=[\,v_1\ \cdots\ v_n\,]$ consists of such a basis and
$Av_j=\lambda_jv_j$, then
\[
 P^{-1}AP=\diag(\lambda_1,\ldots,\lambda_n).
\]
\end{theorem}
\begin{proof}
The equality $AP=PD$ is an equality of columns:
its $j$th column says $Av_j=\lambda_jv_j$.
Invertibility of $P$ means that its columns form a basis.
Multiplying by $P^{-1}$ proves both directions.
\end{proof}

The order of the eigenvalues on the diagonal must match the order of the
eigenvectors in $P$.
The preceding rotation matrix is diagonalisable over $\mathbb C$, but not
over $\mathbb R$.

\begin{example}[A matrix that is not diagonalisable]
For
\[
 A=\begin{pmatrix}1&1\\0&1\end{pmatrix},
 \qquad E_1(A)=\spanop((1,0)^{\mathsf T}).
\]
Indeed, $Av=\lambda v$ first gives
$(1-\lambda)y=0$ and $(1-\lambda)x+y=0$.
A nonzero solution forces $\lambda=1$, and then $y=0$.
Thus there are not enough independent eigenvectors to form a basis.
\end{example}

\begin{lemma}[Independence of distinct eigenspaces]
Let $\lambda_1,\ldots,\lambda_r$ be distinct eigenvalues and
$v_j\in E_{\lambda_j}(A)$.
If $v_1+\cdots+v_r=0$, then every $v_j=0$.
\end{lemma}
\begin{proof}
Use induction on $r$. The case $r=1$ is immediate.
Apply $A-\lambda_r I$ to the given relation:
\[
 \sum_{j=1}^{r-1}(\lambda_j-\lambda_r)v_j=0.
\]
The induction hypothesis implies
$(\lambda_j-\lambda_r)v_j=0$ for $j<r$.
Distinctness makes these scalar factors nonzero, so $v_j=0$ for $j<r$;
the original relation then gives $v_r=0$.
\end{proof}

\begin{corollary}
Eigenvectors with distinct eigenvalues are linearly independent.
An $n\times n$ matrix has at most $n$ distinct eigenvalues.
If it has $n$ distinct eigenvalues in $F$, it is diagonalisable over $F$.
\end{corollary}
\begin{proof}
Apply the lemma to $a_jv_j$ in a proposed linear relation.
Nonzero $v_j$ then force all coefficients to vanish.
There cannot be more than $n$ independent vectors in $F^n$, while
$n$ independent vectors form a basis.
\end{proof}

\begin{theorem}[Dimension criterion]
If $\lambda_1,\ldots,\lambda_r$ are all the eigenvalues in $F$, then
\[
 A\text{ is diagonalisable over }F
 \quad\Longleftrightarrow\quad
 \sum_{j=1}^r\dim E_{\lambda_j}(A)=n.
\]
\end{theorem}
\begin{proof}
Choose a basis in each eigenspace. The preceding lemma shows that their
combined sequence is independent: group a linear relation according to
eigenvalue, then use independence within each eigenspace.
The combined sequence is an eigenvector basis exactly when it has
length $n$. If an eigenvector basis already exists, each of its vectors
belongs to one of these eigenspaces, so their sum is all of $F^n$.
\end{proof}

Distinct eigenvalues are sufficient, but not necessary.
For example, $I_n$ is already diagonal although it has only one
eigenvalue.

\section{Determinants}

The source uses the determinant and its properties without constructing
it from first principles. We follow that scope here.
For a square matrix, $\det A$ is a scalar satisfying, in particular,
\[
 \det I_n=1,\qquad \det(AB)=\det A\,\det B,\qquad
 A\text{ invertible}\ \Longleftrightarrow\ \det A\ne0.
\]
The last condition extends the invertibility theorem in
Chapter~\ref{chap:algebra-subspaces}.
The determinant is not generally additive: $\det(A+B)$ need not equal
$\det A+\det B$.

For small matrices,
\[
 \det\begin{pmatrix}a&b\\c&d\end{pmatrix}=ad-bc,
\]
and expansion along the first row gives
\[
 \det\begin{pmatrix}a&b&c\\d&e&f\\g&h&i\end{pmatrix}
 =a(ei-fh)-b(di-fg)+c(dh-eg).
\]
The signs in this expansion are $+,-,+$.

\subsection{Computing determinants by elimination}

The elementary row operations affect determinants as follows.
\begin{center}
\begin{tabular}{ll}
\toprule
Operation & Effect on determinant\\
\midrule
Swap two rows & Multiply by $-1$\\
Multiply one row by $c$ & Multiply by $c$\\
Add a multiple of one row to another & No change\\
\bottomrule
\end{tabular}
\end{center}
For an upper or lower triangular matrix,
\[
 \boxed{\det A=a_{11}a_{22}\cdots a_{nn}.}
\]
For an upper triangular matrix with nonzero diagonal, eliminate entries
above the diagonal using row additions to obtain a diagonal matrix;
its determinant is the product of its diagonal entries.
If a diagonal entry is zero, such a triangular matrix is singular
(its homogeneous system has a nonzero solution by backward solving),
so its determinant and its diagonal product are both zero.
The lower triangular case follows similarly.

\method{Reduce to triangular form, keeping track of every row swap and
row scaling. Multiply the diagonal entries, then undo the recorded
factors. Using only row additions avoids scaling factors.}

\begin{example}
Using row additions alone,
\[
 \begin{pmatrix}1&2&3\\4&5&6\\7&8&10\end{pmatrix}
 \longrightarrow
 \begin{pmatrix}1&2&3\\0&-3&-6\\0&-6&-11\end{pmatrix}
 \longrightarrow
 \begin{pmatrix}1&2&3\\0&-3&-6\\0&0&1\end{pmatrix}.
\]
Thus the determinant is $-3$.
Replacing the last entry $10$ by $9$ gives a zero row after elimination
and determinant zero.
The cyclic permutation matrix
$\begin{pmatrix}0&1&0\\0&0&1\\1&0&0\end{pmatrix}$ has determinant $1$,
because two row swaps turn it into $I_3$.
\end{example}

\subsection{The cofactor inverse formula}

For reference, the larger inverse formula mentioned in
Chapter~\ref{chap:algebra-matrices} has a compact description.
Delete row $i$ and column $j$ to obtain the minor matrix $A_{ij}$, and
define the cofactor $C_{ij}=(-1)^{i+j}\det A_{ij}$.
The \emph{adjugate} is the transpose of the cofactor matrix:
$\operatorname{adj}(A)=(C_{ji})$.
Cofactor expansion gives
\[
 A\operatorname{adj}(A)=\operatorname{adj}(A)A=(\det A)I_n.
\]
Indeed, a diagonal entry is a determinant expansion; an off-diagonal
entry is the expansion of a matrix with two identical rows or columns,
whose determinant is zero.
Consequently,
\[
 \boxed{A^{-1}=\frac{\operatorname{adj}(A)}{\det A}\quad(\det A\ne0).}
\]
For larger matrices, row reduction is usually more practical.
This cofactor explanation supplements the explicit small-matrix formulae
in the source.

\section{The characteristic polynomial}

\begin{definition}
The \emph{characteristic polynomial} of $A$ is
\[
 c_A(t)=\det(tI_n-A).
\]
It is a monic polynomial of degree $n$.
\end{definition}

The convention $\det(tI-A)$ fixes the leading sign.
Some texts use $\det(A-tI)$, which differs by the factor $(-1)^n$
and has the same roots.

\begin{theorem}
$\lambda\in F$ is an eigenvalue of $A$ if and only if $c_A(\lambda)=0$.
\end{theorem}
\begin{proof}
\[
 c_A(\lambda)=0
 \ \Longleftrightarrow\ \lambda I_n-A\text{ is singular}
 \ \Longleftrightarrow\ N(\lambda I_n-A)\ne\{0_n\}.
\]
A nonzero vector in this kernel satisfies $Av=\lambda v$.
\end{proof}

For $A=\begin{pmatrix}a&b\\c&d\end{pmatrix}$,
\[
 c_A(t)=(t-a)(t-d)-bc=t^2-(a+d)t+(ad-bc).
\]
For a triangular matrix,
$c_A(t)=\prod_{j=1}^n(t-a_{jj})$, so its eigenvalues are precisely its
diagonal entries.

\begin{corollary}[Source Corollary 4.21]
Every complex $n\times n$ matrix with $n\geq1$ has an eigenvalue and a
nonzero eigenvector over $\mathbb C$.
\end{corollary}
\begin{proof}
The fundamental theorem of algebra gives a complex root of $c_A$.
The theorem then gives a nonzero vector in the corresponding kernel.
\end{proof}

This conclusion need not hold over $\mathbb R$: the real rotation example
has characteristic polynomial $t^2+1$.
Existence of a complex eigenvector also does not imply existence of a
whole eigenvector basis.

\subsection{Repeated roots}

The \emph{algebraic multiplicity} of an eigenvalue is its multiplicity as
a root of $c_A(t)$.
Its \emph{geometric multiplicity} is $\dim E_\lambda(A)$.
This terminology helps distinguish repeated roots from the number of
independent eigenvectors they supply.
For $J$, $c_J(t)=t^2(t-3)$: the zero eigenvalue has both multiplicities
two, and the eigenvalue $3$ has both multiplicities one.
For $\begin{pmatrix}1&1\\0&1\end{pmatrix}$,
$c_A(t)=(t-1)^2$ but $\dim E_1(A)=1$.

\section{A complete diagonalisation procedure}

\begin{enumerate}
\item Find $c_A(t)$ and its roots in the chosen field $F$.
\item For each distinct root $\lambda$, row-reduce $A-\lambda I_n$ and
      find a basis of $E_\lambda(A)$.
\item Add the eigenspace dimensions. If the total is less than $n$,
      $A$ is not diagonalisable over $F$.
\item Otherwise place the eigenvectors in the columns of $P$ and their
      eigenvalues, in the same order, on the diagonal of $D$.
\item Verify $AP=PD$. This checks the result without explicitly computing
      $P^{-1}$.
\end{enumerate}

\begin{example}[Three distinct eigenvalues]
For
\[
 A=\begin{pmatrix}1&2&3\\0&4&5\\0&0&6\end{pmatrix},
 \qquad c_A(t)=(t-1)(t-4)(t-6).
\]
Solving the three kernel systems gives, with convenient integer scalings,
\[
 E_1=\spanop\left(\begin{pmatrix}1\\0\\0\end{pmatrix}\right),\quad
 E_4=\spanop\left(\begin{pmatrix}2\\3\\0\end{pmatrix}\right),\quad
 E_6=\spanop\left(\begin{pmatrix}16\\25\\10\end{pmatrix}\right).
\]
Thus
\[
 P=\begin{pmatrix}1&2&16\\0&3&25\\0&0&10\end{pmatrix},\qquad
 D=\diag(1,4,6),\qquad AP=PD.
\]
Rescaling any eigenvector by a nonzero scalar changes $P$ but does not
change its eigenvalue.
\end{example}

\begin{example}[A repeated eigenvalue that causes no problem]
For the all-ones matrix,
\[
 P=\begin{pmatrix}-1&-1&1\\1&0&1\\0&1&1\end{pmatrix},\qquad
 D=\diag(0,0,3).
\]
The first two columns form a basis of $E_0(J)$ and the last spans $E_3(J)$.
They form an eigenvector basis, so $P^{-1}JP=D$.
\end{example}

\begin{proposition}[Only one eigenvalue]
Suppose the only eigenvalue of $A$ in $F$ is $\lambda$.
Then $A$ is diagonalisable over $F$ if and only if $A=\lambda I_n$.
\end{proposition}
\begin{proof}
If $A$ is diagonalisable, its diagonal form must be $\lambda I_n$.
Then $A=P(\lambda I_n)P^{-1}=\lambda I_n$.
The converse holds because $\lambda I_n$ is already diagonal.
\end{proof}

\begin{example}
The matrix
\[
 \begin{pmatrix}
 2&2&3&4\\0&2&5&6\\0&0&2&7\\0&0&0&2
 \end{pmatrix}
\]
has characteristic polynomial $(t-2)^4$.
It is not $2I_4$, so it is not diagonalisable.
\end{example}
